\documentclass[12pt]{article}
\usepackage[spanish]{babel}
\usepackage[utf8]{inputenc}
\usepackage{graphicx}
\usepackage{geometry}
\usepackage{parskip}
\usepackage{setspace}
\usepackage{titlesec}
\usepackage{tabularx}
\usepackage{booktabs}
\usepackage{hyperref}
\usepackage{float}
\usepackage{longtable}
\usepackage{array}
\usepackage{colortbl}
\usepackage{xcolor}
\usepackage{tikz}
\usepackage{subcaption}
\usepackage{quoting}
\usepackage{enumitem}

\geometry{a4paper, margin=2.5cm}
\setlength{\parindent}{0pt}
\renewcommand{\baselinestretch}{1.15}

% Configuración de títulos
\titleformat{\section}{\large\bfseries}{}{0em}{}
\titleformat{\subsection}{\normalsize\bfseries}{}{0em}{}
\titleformat{\subsubsection}{\normalsize\bfseries\itshape}{}{0em}{}

% Colores
\definecolor{high}{rgb}{0.9,0.2,0.2}
\definecolor{medium}{rgb}{0.95,0.7,0.2}
\definecolor{low}{rgb}{0.2,0.7,0.3}

% Entorno para sangría francesa
\newenvironment{hangparas}{\begin{quoting}[hang=true, indent=0.5cm]}{\end{quoting}}

% ---- BIBLIOGRAFÍA UNIFICADA (biblatex APA 7) ----
\usepackage[backend=biber, style=apa, citestyle=apa]{biblatex}
\DeclareLanguageMapping{spanish}{spanish-apa}

\begin{filecontents*}{referencias.bib}
@book{galin2018,
    title = {Software Quality: Concepts and Practice},
    author = {Galin, Daniel},
    year = {2018},
    publisher = {Wiley-IEEE Computer Society Press},
    isbn = {978-1-119-13449-7}
}
@book{pressman2020,
    title = {Software Engineering: A Practitioner's Approach},
    author = {Pressman, Roger S. and Maxim, Bruce R.},
    year = {2020},
    publisher = {McGraw-Hill},
    edition = {9th},
    isbn = {978-1-259-87297-6}
}
@standard{ieee730,
    title = {IEEE Standard for Software Quality Assurance Processes},
    author = {{IEEE Computer Society}},
    year = {2014},
    number = {IEEE Std 730-2014},
    doi = {10.1109/IEEESTD.2014.6835311}
}
@standard{iso25010,
    title = {ISO/IEC 25010:2023 — SQuaRE — Product quality model},
    author = {{ISO/IEC}},
    year = {2023}
}
@standard{iso25023,
    title = {ISO/IEC 25023:2016 — Measurement of product quality},
    author = {{ISO/IEC}},
    year = {2016}
}
@book{oktaba2008,
    title = {Software Process Improvement for Small and Medium Enterprises},
    author = {Oktaba, Hanna and Piattini, Mario},
    year = {2008},
    publisher = {IGI Global}
}
@standard{nmx059,
    title = {NMX-I-059-NYCE-2011 — MoProSoft},
    author = {{NYCE} / {Secretaría de Economía}},
    year = {2011}
}
@standard{cmmi2018,
    title = {CMMI for Development v2.0},
    author = {{CMMI Institute}},
    year = {2018}
}
@standard{iso33000,
    title = {ISO/IEC 33000:2015 — Process assessment},
    author = {{ISO/IEC}},
    year = {2015}
}
@misc{sonarqube,
    title = {SonarQube Documentation},
    author = {{SonarSource}},
    year = {2024},
    url = {https://docs.sonarqube.org/}
}
@misc{istqb2023,
    title = {ISTQB Foundation Level Syllabus v4.0},
    author = {{ISTQB}},
    year = {2023},
    url = {https://www.istqb.org/}
}
@book{forsgren2018,
    title = {Accelerate: The Science of Lean Software and DevOps},
    author = {Forsgren, Nicole and Humble, Jez and Kim, Gene},
    year = {2018},
    publisher = {IT Revolution}
}
@book{piattini2019,
    title = {Calidad de Sistemas de Información},
    author = {Piattini, Mario G.},
    year = {2019},
    publisher = {RA-MA},
    edition = {5th}
}
\end{filecontents*}
\addbibresource{referencias.bib}

\begin{document}

% ==================== PORTADA (Olivia) ====================
\begin{titlepage}
    \thispagestyle{empty}
    \begin{center}
        \textsc{\Large {\textbf{UNIVERSIDAD POLITÉCNICA DE AGUASCALIENTES}}}\\
        {\Large \uppercase{Ingeniería en Tecnologías de la Información e Innovación Digital} \\ }

         \begin{center}
            % APARTADO PARA IMAGEN 1: Logo UPA
            \includegraphics[width=0.45\textwidth]{imagenes/LOGO UPA.png}
            % \hspace{0.5cm}
            % APARTADO PARA IMAGEN 2: Logo TIID
             \includegraphics[width=0.35\textwidth]{imagenes/TIID UPA.png}\
        \end{center}

        \vspace{0.1cm}
        \rule{\linewidth}{0.3mm} \\
        { \Huge \bfseries Tarea 03: Aplicación de estándares\\[0.2cm]
          \large (ISO/IEC 25010, IEEE 730, CMMI, MoProSoft) }
        \rule{\linewidth}{0.3mm} \\

        \vspace{0.1cm}

        {\large \textbf{ASIGNATURA}}\\
        Estándares y Métricas para el Desarrollo de Software\\

        \vspace{0.4cm}

        {\large \textbf{DOCENTE}}\\
        Dr(c). Raul Alejandro Velasquez Ortiz \\

        \vspace{0.4cm}

        {\large \textbf{PROYECTO INTEGRADOR}}\\
        Outline - Wiki interna estilo Notion (Open Source maduro) \\

        \vspace{0.4cm}

        {\large \textbf{EQUIPO No. 4}}\\
        \vspace{0.4cm}
        \textbf{CHAIREZ ALBA OLIVIA LIZBETH - UP240594} \\
        \textbf{ESTRADA RAMÍREZ NOÉ EZEQUIEL - UP240770} \\
        \textbf{FACIO PALOS OMAR - UP240470} \\
        \textbf{JAUREGUI DE LA RIVA ALAN RICARDO - UP240960} \\
        \textbf{RUBIO VIRAMONTES BRIAN ODIN - UP250004}

        \vspace{0.4cm}

        {\large \textbf{GRUPO}}\\
        TIID05B

        \vfill
        {\large \textbf{Repositorio Git:} \url{https://github.com/UP240594/outline-sqa-docs}}

        {\large Aguascalientes, México\\
        18 de mayo de 2026}
    \end{center}
\end{titlepage}

% ==================== ÍNDICE ====================
\newpage
\tableofcontents
\newpage

% ==================== 1. INTRODUCCIÓN (Olivia) ====================
\section{Introducción}

\subsection{Contexto del proyecto}
El presente documento corresponde a la Tarea 03 de la asignatura Estándares y Métricas para el Desarrollo de Software, en la cual se aplican los conceptos teóricos de calidad de software al proyecto integrador del equipo: \textbf{Outline}, una aplicación web de tipo wiki interna con funcionalidades similares a Notion, de código abierto y con una comunidad activa.

Outline fue seleccionado del Catálogo de 30 alternativas (Semana 1) por su naturaleza madura (más de 5,000 estrellas en GitHub, licencia MIT) y por ser un caso realista para evaluar atributos de calidad como usabilidad, seguridad y fiabilidad. El proyecto será analizado desde la perspectiva de un equipo de desarrollo que busca adoptar, desplegar y posiblemente contribuir a este software en un entorno corporativo.

\subsection{Objetivos de calidad}
\begin{itemize}
    \item Aplicar el modelo de calidad de producto ISO/IEC 25010:2023 para priorizar las características más relevantes para Outline.
    \item Definir KPIs cuantificables basados en ISO/IEC 25023.
    \item Estructurar un Plan SQA preliminar siguiendo IEEE 730-2014.
    \item Comparar modelos de madurez y recomendar el más adecuado.
    \item Identificar riesgos de calidad y proponer mitigaciones.
\end{itemize}

\subsection{Estructura del documento}
El documento se organiza en seis secciones principales: descripción del proyecto, análisis ISO 25010, KPIs, plan SQA, modelo de madurez y matriz de riesgos. Finaliza con conclusión, aportes individuales y referencias.

% ==================== 2.1 DESCRIPCIÓN DEL PROYECTO (trabajo de todos + diagrama añadido) ====================
\section{Descripción del proyecto integrador}

\subsection{Nombre y tipo}
\textbf{Proyecto:} Outline - Wiki interna estilo Notion \\
\textbf{Tipo:} Web (SaaS con opción self-hosted)

\subsection{Contexto de uso}
Outline está diseñado para equipos que necesitan un repositorio centralizado de conocimiento interno. Ofrece despliegue self-hosted gratuito, integración OAuth y editor colaborativo en tiempo real.

\subsection{Usuarios objetivo}
\begin{table}[H]
\centering
\begin{tabularx}{\textwidth}{|l|X|}
\hline
\textbf{Rol} & \textbf{Necesidad principal} \\ \hline
Editor de contenido & Crear y organizar documentos técnicos. \\ \hline
Lector/Consultor & Buscar información rápidamente. \\ \hline
Administrador & Gestionar permisos e invitaciones. \\ \hline
Administrador de sistemas & Desplegar y mantener la instancia. \\ \hline
\end{tabularx}
\caption{Usuarios objetivo de Outline}
\end{table}

\subsection{Stack tecnológico}
\begin{itemize}
    \item Frontend: React 18+, TypeScript, MobX, Slate.js
    \item Backend: Node.js + Express, Sequelize
    \item Bases de datos: PostgreSQL, Redis
    \item Infraestructura: Docker Compose, AWS S3/MinIO
\end{itemize}

\subsection{Alcance en historias de usuario}
\begin{enumerate}
    \item HU-01: Creación de documento
    \item HU-02: Edición colaborativa en tiempo real
    \item HU-03: Gestión de permisos y equipos
    \item HU-04: Búsqueda y navegación
    \item HU-05: Historial de versiones
\end{enumerate}

\subsection{Diagrama de alto nivel (arquitectura)}
\begin{figure}[H]
\centering
\begin{tikzpicture}[node distance=1.5cm, auto, block/.style={rectangle, draw, thick, text width=2.5cm, align=center, minimum height=1cm}]
    \node[block] (user) {Usuario (Navegador)};
    \node[block, right of=user, node distance=4cm] (frontend) {Frontend React + TypeScript};
    \node[block, right of=frontend, node distance=4cm] (api) {API Node.js + Express};
    \node[block, below of=api, node distance=2cm] (db) {PostgreSQL};
    \node[block, left of=db, node distance=3cm] (redis) {Redis (caché)};
    \draw[->, thick] (user) -- (frontend);
    \draw[->, thick] (frontend) -- (api);
    \draw[->, thick] (api) -- (db);
    \draw[<->, thick] (api) -- (redis);
\end{tikzpicture}
\caption{Diagrama de alto nivel de la arquitectura de Outline. Fuente: Elaboración propia.}
\label{fig:diagrama}
\end{figure}

% ==================== 2.2 ISO 25010 (Noé) ====================
\section{Análisis de calidad con ISO/IEC 25010}

El estándar internacional ISO/IEC 25010:2023 proporciona el marco estructural para diagnosticar y asegurar la calidad del producto de software. En el contexto de \textbf{Outline}, concebido como una plataforma de gestión del conocimiento de alto rendimiento y estilo wiki, los atributos de calidad han sido mapeados y priorizados de acuerdo con su impacto crítico en entornos de productividad empresarial distribuidos.

\subsection{Priorización de Características de Calidad del Producto}

\begin{itemize}
    \item \textbf{Adecuación Funcional [Prioridad: ALTA]} \\
    Outline debe garantizar que sus funcionalidades de edición colaborativa en tiempo real, renderizado de Markdown y control de versiones cumplan con exactitud las expectativas del usuario corporativo.
    
    \item \textbf{Eficiencia de Desempeño [Prioridad: ALTA]} \\
    Los tiempos de carga de documentos extensos y la latencia del motor de búsqueda deben ser mínimos, asegurando un uso fluido.
    
    \item \textbf{Usabilidad [Prioridad: ALTA]} \\
    La tasa de adopción está ligada a la simplicidad de la experiencia de usuario; una interfaz intuitiva garantiza que personal no técnico colabore sin capacitación compleja.
    
    \item \textbf{Fiabilidad [Prioridad: ALTA]} \\
    La tolerancia a fallos de sincronización y los mecanismos de respaldo automático son indispensables. La pérdida de documentación interna interrumpe los flujos operativos.
    
    \item \textbf{Seguridad [Prioridad: ALTA]} \\
    Es obligatorio implementar control de acceso RBAC, soporte SSO/SAML y cifrado de datos en tránsito.
    
    \item \textbf{Compatibilidad [Prioridad: MEDIA]} \\
    La wiki requiere integración fluida con herramientas externas (Slack, Google Drive, Figma).
    
    \item \textbf{Mantenibilidad [Prioridad: MEDIA]} \\
    El diseño con NodeJS y React facilita la modularidad. Al ser open source, el equipo debe asegurar legibilidad para integraciones a medida.
    
    \item \textbf{Portabilidad [Prioridad: MEDIA]} \\
    La instalación se mitiga mediante el uso de contenedores Docker y Kubernetes.
\end{itemize}

\begin{figure}[H]
    \centering
    \begin{tikzpicture}[scale=0.85]
        \draw[very thin, gray!30] (0,0) grid (10,8.5);
        \draw[thick, ->] (0,0) -- (10.5,0) node[anchor=north east] {Nivel de Prioridad};
        \draw[thick, ->] (0,0) -- (0,9) node[anchor=south east] {Características};
        \node at (2.5,-0.4) {\footnotesize Baja};
        \node at (6,-0.4) {\footnotesize Media};
        \node at (9,-0.4) {\footnotesize Alta};
        \fill[fill=low] (0,0.5) rectangle (6,1.1); \node[anchor=west] at (6,0.8) {\scriptsize Portabilidad};
        \fill[fill=low] (0,1.5) rectangle (6,2.1); \node[anchor=west] at (6,1.8) {\scriptsize Mantenibilidad};
        \fill[fill=low] (0,2.5) rectangle (6,3.1); \node[anchor=west] at (6,2.8) {\scriptsize Compatibilidad};
        \fill[fill=high] (0,3.5) rectangle (9,4.1); \node[anchor=west] at (9,3.8) {\scriptsize Seguridad};
        \fill[fill=high] (0,4.5) rectangle (9,5.1); \node[anchor=west] at (9,4.8) {\scriptsize Fiabilidad};
        \fill[fill=high] (0,5.5) rectangle (9,6.1); \node[anchor=west] at (9,5.8) {\scriptsize Usabilidad};
        \fill[fill=high] (0,6.5) rectangle (9,7.1); \node[anchor=west] at (9,6.8) {\scriptsize Eficiencia};
        \fill[fill=high] (0,7.5) rectangle (9,8.1); \node[anchor=west] at (9,7.8) {\scriptsize Adecuación Funcional};
    \end{tikzpicture}
    \caption{Visualización gráfica del perfil de prioridades de calidad para Outline.}
\end{figure}

\subsection{Análisis de Calidad en Uso}
Los dos componentes con mayor impacto son:
\begin{enumerate}
    \item \textbf{Efectividad:} Capacidad del software para permitir que los usuarios alcancen metas con precisión. En Outline, se materializa en el guardado síncrono y la precisión del buscador.
    \item \textbf{Satisfacción:} Respuesta cognitiva y emocional del usuario. El diseño minimalista y los accesos rápidos elevan la satisfacción.
\end{enumerate}

% ==================== 2.3 KPIs (Omar) ====================
\section{KPIs medibles según ISO/IEC 25023}

Con el objetivo de cuantificar objetivamente las características de calidad priorizadas, el equipo ha definido 7 KPIs basados en la norma \cite{iso25023}.

\begin{table}[H]
\centering
\begin{tabularx}{\textwidth}{|>{\columncolor[gray]{0.92}}p{3.2cm}|>{\centering\arraybackslash}p{2.5cm}|p{3.5cm}|>{\centering\arraybackslash}p{2.5cm}|p{2.2cm}|}
\hline
\rowcolor[gray]{0.85}
\textbf{Nombre del KPI} & \textbf{Caract. ISO 25010} & \textbf{Fórmula} & \textbf{Umbral} & \textbf{Herramienta} \\ \hline
Tasa de éxito de creación de documentos & Adecuación funcional & \(\displaystyle \frac{\text{Docs creados exitosamente}}{\text{Intentos de creación}} \times 100\) & \(\geq 99.5\%\) & Logs de API + Grafana \\ \hline
Disponibilidad del servicio & Fiabilidad & \(\displaystyle \frac{T_{\text{up}}}{T_{\text{up}} + T_{\text{inactivo}}}\) & \(\geq 99.9\%\) & Uptime Kuma \\ \hline
Tiempo de carga del dashboard & Eficiencia (rendimiento) & Percentil 95 (P95) en ms & \(< 1.5\) s & Lighthouse CI \\ \hline
Tiempo de respuesta de búsqueda & Eficiencia & P95 de latencia del endpoint & \(< 500\) ms & Prometheus + Grafana \\ \hline
Cobertura de pruebas de seguridad & Seguridad & \(\displaystyle \frac{\text{Pruebas ejecutadas}}{\text{Total casos}} \times 100\) & \(\geq 80\%\) & OWASP ZAP + SonarQube \\ \hline
Tasa de retención de usuarios & Calidad en uso & \(\displaystyle \frac{\text{Usuarios activos 30d}}{\text{Usuarios nuevos mes}} \times 100\) & \(\geq 70\%\) & Analytics propio \\ \hline
Tiempo para primera edición & Usabilidad & Tiempo desde registro hasta primera edición & \(< 3\) min & Session replay \\ \hline
\end{tabularx}
\caption{KPIs de Calidad del Producto}
\end{table}

\subsection*{Nota metodológica}
\begin{enumerate}
    \item \textbf{Línea base:} Establecer medición inicial después del primer sprint \cite{galin2018}.
    \item \textbf{Automatización:} KPIs de rendimiento y seguridad deben medirse en cada commit mediante CI/CD \cite{forsgren2018}.
    \item \textbf{Justificación de umbrales:} Basada en benchmarks de la industria y recomendaciones de Galin.
    \item \textbf{Percentil P95:} Refleja mejor la experiencia real del usuario que el promedio \cite{pressman2020}.
\end{enumerate}

% ==================== 2.4 PLAN SQA (Alan, corregido) ====================
\section{Plan SQA preliminar según IEEE 730}

\subsection{§1 Purpose (Propósito y alcance)}
Plan de Aseguramiento de la Calidad del Software conforme al estándar IEEE 730-2014 \cite{ieee730} para el proyecto \textbf{Outline}. Establece actividades, responsabilidades y métricas para garantizar el cumplimiento de requisitos funcionales y no funcionales.

\textbf{Alcance:} Aplica a todos los entregables del proyecto Outline: código fuente, documentación técnica, manuales de usuario, informes de pruebas. Cubre las fases de análisis, diseño, implementación, pruebas y despliegue.

\subsection{§3 Management (Roles y responsabilidades)}
\begin{table}[H]
\centering
\caption{Roles y responsabilidades en el plan SQA (adaptado al equipo)}
\begin{tabular}{|l|p{8cm}|}
\hline
\textbf{Rol} & \textbf{Responsabilidades} \\ \hline
Gerente de Proyecto (Noé E. Estrada) & Aprobar el plan SQA, asignar recursos, supervisar hitos de calidad. \\ \hline
Líder de Calidad / QA Tester (Omar Facio) & Ejecutar auditorías internas, medir KPIs, reportar no conformidades. \\ \hline
Desarrolladores (Olivia Chairez, Odin Rubio) & Seguir estándares de codificación, realizar pruebas unitarias. \\ \hline
Documentador (Alan Jauregui) & Mantener actualizada la documentación de calidad. \\ \hline
Cliente / Usuario clave & Validar criterios de aceptación, reportar problemas en fase de pruebas. \\ \hline
\end{tabular}
\end{table}

\subsection{§5 Standards, practices, metrics}
\textbf{Estándares:} PEP 8 (Python), ESLint (JavaScript), \LaTeX\ para documentación, ISTQB v4.0 para pruebas, Git Flow. \\
\textbf{Prácticas:} Commits con ID de tarea, pull requests con revisión y análisis estático, reuniones semanales con revisión de KPIs. \\
\textbf{Métricas:} Disponibilidad diaria, cobertura de pruebas por sprint, tiempo de respuesta semanal, etc.

\subsection{§7 Test (Estrategia de pruebas)}
\begin{enumerate}
    \item Unitarias: JUnit/Pytest, cobertura mínima 80\%.
    \item Integración: Postman/Newman en cada CI.
    \item Sistema: Selenium, JMeter.
    \item Aceptación: Con cliente en preproducción.
\end{enumerate}
Automatización con GitHub Actions.

\subsection{§8 Problem reporting (SLA por severidad)}
\begin{table}[H]
\centering
\caption{SLA para reporte de problemas}
\begin{tabular}{|l|l|l|l|}
\hline
\textbf{Severidad} & \textbf{Descripción} & \textbf{Tiempo respuesta} & \textbf{Solución} \\ \hline
Crítica (S1) & Pérdida de datos, caída total & 2 horas & 24 horas \\ \hline
Alta (S2) & Funcionalidad principal afectada & 8 horas & 3 días hábiles \\ \hline
Media (S3) & Funcionalidad secundaria afectada & 2 días hábiles & 7 días hábiles \\ \hline
Baja (S4) & Error estético o mejora menor & 5 días hábiles & Siguiente sprint \\ \hline
\end{tabular}
\end{table}
Reporte mediante GitHub Issues.

% ==================== 2.5 MODELO DE MADUREZ (Odin) ====================
\section{Recomendación de modelo de madurez (CMMI / MoProSoft / ISO 33000)}

\subsection{Análisis comparativo de modelos}
\begin{table}[H]
\centering
\small
\begin{tabularx}{\textwidth}{|l|X|X|X|}
\hline
\textbf{Criterio} & \textbf{CMMI-DEV v2.0} & \textbf{MoProSoft} & \textbf{ISO/IEC 33000} \\ \hline
Origen & EE.UU. & México & Internacional \\ \hline
Enfoque & Grandes contratistas & PYMEs mexicanas & Evaluación de capacidad \\ \hline
Niveles & 5 & 3 & 6 \\ \hline
Costo & Alto & Bajo & Medio \\ \hline
\end{tabularx}
\caption{Comparativa de modelos de madurez}
\end{table}

\subsection{Contexto específico}
Equipo de 5 estudiantes, PYME micro, sector software de productividad, mercado México + internacional.

\subsection{Recomendación: MoProSoft (con elementos de ISO 33000)}
\begin{itemize}
    \item Alineación con tamaño del equipo (diseñado para PYMEs mexicanas) \cite{oktaba2008}.
    \item Costo documental cero (norma NMX-I-059 pública).
    \item Estructura simplificada (Alta Dirección, Gerencia, Operación).
    \item Relevancia para el mercado mexicano.
    \item Complemento con ISO 33000 para medición de capacidad.
\end{itemize}

\subsection{Implementación propuesta}
\begin{itemize}
    \item Alta Dirección: seguimiento de KPIs y asignación de roles.
    \item Gerencia: gestión de requisitos y planificación.
    \item Operación: estándares de codificación y estrategia de pruebas.
\end{itemize}

% ==================== 2.6 MATRIZ DE RIESGOS (Odin) ====================
\section{Matriz de riesgos de calidad}

Se identifican 7 riesgos específicos del proyecto Outline.

\begin{table}[H]
\centering
\footnotesize
\renewcommand{\arraystretch}{1.3}
\begin{tabularx}{\textwidth}{|l|X|c|c|X|}
\hline
\textbf{Riesgo} & \textbf{Característica afectada} & \textbf{Prob.} & \textbf{Impacto} & \textbf{Mitigación} \\ \hline
R1: Exposición de secretos en logs & Seguridad & M & A & git-secrets, rotar secretos \\ \hline
R2: Conflictos en editor colaborativo & Fiabilidad & M & A & Pruebas de estrés, backups cada 5 min \\ \hline
R3: Curva de aprendizaje alta & Usabilidad & A & M & Tutorial interactivo, encuestas \\ \hline
R4: Código extenso (50k+ líneas) & Mantenibilidad & A & M & Documentación con Typedoc \\ \hline
R5: Vulnerabilidades en npm & Seguridad & A & A & npm audit semanal, Snyk \\ \hline
R6: Rendimiento en búsquedas & Eficiencia & M & A & Indexar Elasticsearch, pruebas k6 \\ \hline
R7: Falta de participación & Adecuación funcional & A & M & Reuniones semanales, checklist en PR \\ \hline
\end{tabularx}
\caption{Matriz de riesgos de calidad}
\end{table}

\textbf{Leyenda:} A = Alta, M = Media, B = Baja

\subsection{Plan de contingencia}
\begin{itemize}
    \item \textbf{R5:} Vulnerabilidad crítica sin parche → bloqueador de release, buscar workaround.
    \item \textbf{R2:} Pérdida de datos → shadow log que registra todas las operaciones.
\end{itemize}

% ==================== CONCLUSIÓN (Olivia) ====================
\newpage
\section{Conclusión integral}

La aplicación de los estándares ISO/IEC 25010, IEEE 730, CMMI y MoProSoft al proyecto Outline ha permitido al equipo descubrir varios aprendizajes clave:

\paragraph{Hallazgo 1: La priorización es esencial.} No todas las características de calidad tienen el mismo peso. Para una wiki interna, usabilidad y seguridad dominan sobre portabilidad.

\paragraph{Hallazgo 2: Los KPIs deben ser realistas.} Definir umbrales como disponibilidad 99,9\% es alcanzable; 99,999\% no lo es para un proyecto open-source.

\paragraph{Hallazgo 3: IEEE 730 es estructurado pero adaptable.} Adaptarlo a un equipo de 5 personas requirió recortar secciones no aplicables.

\paragraph{Hallazgo 4: MoProSoft es relevante y vigente.} Sigue siendo pertinente para PYMEs locales, con principios de baja burocracia.

\paragraph{Conexión con próximos parciales:} Los KPIs definidos serán la base para las pruebas del Parcial 2, y la matriz de riesgos alimentará el informe técnico del Parcial 3.

% ==================== APORTES INDIVIDUALES ====================
\newpage
\section{Aportes individuales}

\paragraph{Chairez Alba Olivia Lizbeth (UP240594) - Coordinadora:}
Redacté la portada, introducción, descripción del proyecto (sección 2.1) y conclusión. Unifiqué y armonicé todas las partes de mis compañeros en un solo documento LaTeX, corregí errores de compilación y formateé el documento final. Gestioné el repositorio Git y aseguré la consistencia general. Añadí el diagrama de alto nivel requerido.

\paragraph{Estrada Ramírez Noé Ezequiel (UP240770) - Líder:}
Desarrollé la sección 2.2 completa (análisis ISO 25010), incluyendo la priorización de las 8 características, justificaciones específicas para Outline y el análisis de calidad en uso (efectividad y satisfacción). Además, elaboré el gráfico de barras con TikZ.

\paragraph{Facio Palos Omar (UP240470) - QA Tester:}
Elaboré la sección 2.3 (KPIs según ISO 25023), definiendo 7 indicadores con fórmulas, umbrales numéricos y herramientas de medición. Redacté la nota metodológica justificando el uso del percentil P95 y la automatización en CI/CD.

\paragraph{Jauregui De La Riva Alan Ricardo (UP240960) - Documentador:}
Redacté la sección 2.4 (Plan SQA según IEEE 730), incluyendo propósito, gestión, estándares, estrategia de pruebas y SLA por severidad. Adapté el estándar al contexto del equipo y corregí el nombre del proyecto a Outline.

\paragraph{Rubio Viramontes Brian Odin (UP250004) - Developer:}
Desarrollé la sección 2.5 (recomendación de modelo de madurez) con análisis comparativo CMMI/MoProSoft/ISO 33000, justificando MoProSoft como la opción óptima. También elaboré la sección 2.6 (matriz de riesgos) con 7 riesgos identificados, probabilidad, impacto y estrategias de mitigación.

% ==================== REFERENCIAS ====================
\newpage
\printbibliography[title={Referencias}]

\end{document}